\section{Proofs for Optional Stopping and Random Walks}
\label{sec:annex-p15-optional-stopping-and-random-walks}

\begin{annexitem}{annex:p15.os-ui-stopped}{thm:p15-os-ui-stopped}{The source positive-part and tail argument}
\begin{proof}
Uniform integrability implies \(\sup_n\mathbb E|X_n|<\infty\).
The submartingale convergence theorem therefore gives the finite
integrable limit \(X_\infty\), and Vitali's theorem gives
\(X_n\to X_\infty\) in \(L^1\).
Stopping preserves the submartingale property.
Moreover \(X_n^+\) is a submartingale and bounded optional
sampling gives
\[
 \mathbb E(X_{T\wedge n})^+\leq\mathbb EX_n^+.
 \]
The right side is bounded uniformly in \(n\).
The submartingale convergence theorem applied to the stopped
process gives an integrable finite limit. On \(\{T<\infty\}\)
the process eventually equals \(X_T\); on \(\{T=\infty\}\)
it equals \(X_n\) and tends to \(X_\infty\).
Thus its limit is precisely the specified \(X_T\).

The source tail splitting is now legitimate:
\[
 \begin{split}
 \mathbb E(|X_{T\wedge n}|;|X_{T\wedge n}|>K)
 ={}&\mathbb E(|X_T|;|X_T|>K,T\leq n)\\
 &+\mathbb E(|X_n|;|X_n|>K,T>n).
 \end{split}
 \]
The first term is bounded by the tail of one integrable variable,
and the second by the uniformly integrable original family's tail.
They tend to zero uniformly in \(n\), proving uniform
integrability of the stopped process. Vitali gives its \(L^1\) limit.
Finally bounded optional sampling gives
\(\mathbb EX_0\leq\mathbb EX_{T\wedge n}\leq\mathbb EX_n\).
Pass to the two \(L^1\) limits. Applying the argument to \(-X\)
gives the supermartingale version; applying both inequalities
gives equality for martingales.
\end{proof}
\end{annexitem}

\begin{annexitem}{annex:p15.os-residual-ui}{thm:p15-os-residual-ui}{The residual family supplies exactly the second tail}
\begin{proof}
Let \(R_n=X_n\mathbf1_{\{T>n\}}\).
The first term in the preceding tail decomposition is bounded
by \(\mathbb E(|X_T|;|X_T|>K)\), and the second is exactly
\(\mathbb E(|R_n|;|R_n|>K)\) for \(K>0\).
The two assumptions therefore make the stopped family uniformly
integrable. At every finite stopping outcome it stabilizes;
at infinite stopping outcomes convergence is the stated terminal
hypothesis. Thus \(X_{T\wedge n}\to X_T\) almost surely,
and Vitali gives \(L^1\) convergence.
The finite-time sampling inequality passes to the limit.
The same proof with the opposite sign proves the remaining versions.
\end{proof}
\end{annexitem}

\begin{annexitem}{annex:p15.os-conditional-sampling}{thm:p15-os-conditional-sampling}{Testing on every stopped-information event}
\begin{proof}
Write \(Z_n=Y_{M\wedge n}\).
It is a uniformly integrable submartingale, so its finite
integrable terminal value is \(Y_M\).
Theorem~\ref{thm:p15-os-ui-stopped} applied to \(Z\) and the
stopping time \(L\) gives an integrable sampled value \(Y_L\).
Indeed on \(\{L<\infty\}\), \(L\leq M\) gives
\(Z_L=Y_L\), and on \(\{L=\infty\}\) both sampled values
are the same limit \(Y_M\).
The stopped value is \(\mathcal F_L\)-measurable: its intersection
with \(\{L\leq n\}\) is a finite union of the measurable
events \(\{L=j,Y_j\in B\}\), for any Borel \(B\); its
ambient measurability follows from the countable construction
and the measurable terminal limit.

Fix \(A\in\mathcal F_L\) and put \(A_k=A\cap\{L\leq k\}\).
Then \(A_k\in\mathcal F_{L\wedge k}\).
For \(n\geq k\), bounded conditional sampling of \(Z\) gives
\[
 \mathbb E(Y_L;A_k)
   =\mathbb E(Z_{L\wedge k};A_k)\leq\mathbb E(Z_n;A_k).
 \]
First let \(n\to\infty\), using \(Z_n\to Y_M\) in \(L^1\).
Then let \(k\to\infty\), using integrability of both sampled
values, to obtain the inequality on \(A\cap\{L<\infty\}\).
On \(A\cap\{L=\infty\}\) the two values agree.
Hence \(\mathbb E(Y_L;A)\leq\mathbb E(Y_M;A)\) for every
\(A\in\mathcal F_L\). The defining event criterion for conditional
expectation gives the claimed conditional inequality.
Take \(A=\Omega\) for its mean version, and change signs for
the other martingale classes.
\end{proof}
\end{annexitem}

\begin{annexitem}{annex:p15.os-nonnegative}{thm:p15-os-nonnegative}{Fatou sampling and the infinite-history crossing bound}
\begin{proof}
The nonnegative supermartingale convergence theorem gives
the finite integrable limit \(X_\infty\).
For every stopping time \(T\leq\infty\),
\(X_{T\wedge n}\to X_T\) almost surely, and bounded sampling
and Fatou give
\[
 \mathbb EX_T\leq\liminf_n\mathbb EX_{T\wedge n}
                                               \leq\mathbb EX_0.
 \]
In particular all sampled values are integrable.
They are measurable for their stopped sigma-algebras by the
same countable-piece construction used in the preceding proof.
For \(L\leq M\), \(A\in\mathcal F_L\), and
\(A_k=A\cap\{L\leq k\}\), the bounded times
\(L\wedge k\leq M\wedge n\), \(n\geq k\), give
\[
 \mathbb E(X_{M\wedge n};A_k)\leq\mathbb E(X_L;A_k).
 \]
Fatou applied to the nonnegative integrand on the left gives
\(\mathbb E(X_M;A_k)\leq\mathbb E(X_L;A_k)\).
Let \(k\to\infty\); on \(\{L=\infty\}\), also \(M=\infty\)
and both values equal \(X_\infty\).
The conditional event criterion proves the asserted inequality.

For \(c>0\), let \(T_c=\inf\{n:X_n>c\}\), a stopping time
since the search is over adapted coordinates.
Nonnegativity and the first sampling bound give
\[
 c\,\mathbb P(T_c<\infty)\leq\mathbb EX_{T_c}\leq\mathbb EX_0.
 \]
This proves the strict-supremum bound.
For \(0<c<\lambda\), the event \(\{\sup_nX_n\geq\lambda\}\)
is contained in \(\{T_c<\infty\}\); let \(c\uparrow\lambda\).
This handles a supremum equal to \(\lambda\) that need not
be attained at any finite time.
\end{proof}
\end{annexitem}

\begin{annexitem}{annex:p15.os-integrable-time}{thm:p15-os-integrable-time}{An integrable total-variation majorant}
\begin{proof}
Finite expectation makes \(T\) almost surely finite.
The stopped-sum identity gives pointwise
\[
 |X_{T\wedge n}|
 \leq D:=|X_0|+
         \sum_{m=0}^\infty|X_{m+1}-X_m|\mathbf1_{\{T>m\}}.
 \]
Since \(\{T>m\}\in\mathcal F_m\), the defining conditional
expectation identity and the given bound imply
\[
 \mathbb E(|X_{m+1}-X_m|;T>m)\leq B\,\mathbb P(T>m).
 \]
Tonelli and the pointwise identity
\(T=\sum_{m\geq0}\mathbf1_{\{T>m\}}\) give
\(\mathbb ED\leq\mathbb E|X_0|+B\mathbb ET<\infty\).
An integrable majorant implies uniform integrability, and
dominated convergence gives \(X_{T\wedge n}\to X_T\)
in \(L^1\), with the stated absolute-mean bound.
The bounded optional inequality now passes to the limit.
\end{proof}
\end{annexitem}

\begin{annexitem}{annex:p15.os-wald-first}{thm:p15-os-wald-first}{The linear martingale at an integrable time}
\begin{proof}
Independence makes \(M_n=S_n-n\mu\) a martingale of initial
value \(x\), and
\[
 \mathbb E(|M_{n+1}-M_n|\mid\mathcal F_n)
                    =\mathbb E|\xi_1-\mu|<\infty.
 \]
The integrable-time theorem applies, giving \(M_T\in L^1\)
and \(\mathbb EM_T=x\).
Because \(\mu T\) is integrable, \(S_T=M_T+\mu T\) is
integrable and the asserted identity follows.
Alternatively the absolute stopped-increment sum has expectation
\[
 \mathbb E\sum_{j\geq1}|\xi_j|\mathbf1_{\{T\geq j\}}
 =\mathbb E|\xi_1|\sum_{j\geq1}\mathbb P(T\geq j)
 =\mathbb E|\xi_1|\,\mathbb ET,
 \]
which verifies directly the integrability used in that identity.
\end{proof}
\end{annexitem}

\begin{annexitem}{annex:p15.os-wald-second}{thm:p15-os-wald-second}{Stopped orthogonality before the second identity}
\begin{proof}
For every finite horizon,
\[
 S_{T\wedge n}=\sum_{j=1}^n\xi_j\mathbf1_{\{T\geq j\}}.
 \]
The indicator is \(\mathcal F_{j-1}\)-measurable.
The terms have zero means and squared means
\(\sigma^2\mathbb P(T\geq j)\).
Distinct terms are orthogonal: for \(i<j\), condition their
integrable product on \(\mathcal F_{j-1}\), using
\(\mathbb E(\xi_j\mid\mathcal F_{j-1})=0\).
The products are integrable by Cauchy--Schwarz.
Thus
\[
 \mathbb ES_{T\wedge n}^2=\sigma^2\mathbb E(T\wedge n),
 \qquad
 \mathbb E(S_{T\wedge m}-S_{T\wedge n})^2
   =\sigma^2\mathbb E((T\wedge m)-(T\wedge n))\quad(m\geq n).
 \]
When \(\mathbb ET<\infty\), the time is finite almost surely
and \(S_{T\wedge m}\to S_T\) pointwise there.
Fatou first gives \(\mathbb ES_T^2\leq\sigma^2\mathbb ET<\infty\),
then, for fixed \(n\),
\[
 \mathbb E(S_T-S_{T\wedge n})^2
 \leq\sigma^2\mathbb E(T-T\wedge n)\to0.
 \]
This establishes \(L^2\) convergence without assuming its limit's
second moment. Cauchy--Schwarz applied to
\(|A^2-B^2|\leq|A-B|(|A|+|B|)\) now gives convergence
of the squares in \(L^1\), and hence
\(\mathbb ES_T^2=\sigma^2\mathbb ET\).

The source quadratic-martingale route is valid with precisely
these premises: \(Q_n=S_n^2-n\sigma^2\) is a martingale,
since independence and centering give
\(\mathbb E(S_{n+1}^2\mid\mathcal F_n)=S_n^2+\sigma^2\),
and \(Q_{T\wedge n}\) converges in \(L^1\) because both its
square term and its time term do so. It is uniformly integrable.
Its residual \(Q_n\mathbf1_{\{T>n\}}\), obtained by restricting
\(Q_{T\wedge n}\), is uniformly integrable as well.
The residual-tail stopping theorem therefore yields
\(\mathbb EQ_T=\mathbb EQ_0=0\), the same identity.

For the exit-time assertion, if \(\mathbb ET=\infty\) the
extended inequality is immediate. Otherwise \(T<\infty\)
almost surely and \(|S_T|>a\), so
\(\sigma^2\mathbb ET=\mathbb ES_T^2>a^2\).
No claim about a finite expectation was needed in the infinite case.
When \(\sigma^2=0\), each nonnegative squared increment has
expectation zero. A countable intersection gives all increments
zero almost surely, proving the degenerate assertions directly.
\end{proof}
\end{annexitem}

\begin{annexitem}{annex:p15.os-independent-restart}{thm:p15-os-independent-restart}{A stopped-index product calculation}
\begin{proof}
Let \(A\in\mathcal F_T\), and let \(g\) be any bounded Borel
function of \(k\) increments, \(k\geq1\).
The event \(A\cap\{T=n\}\) belongs to \(\mathcal F_n\).
Independence of the fixed future block and identical laws give
\[
 \begin{split}
 &\mathbb E\bigl[\mathbf1_A
       g(\xi_{T+1},\ldots,\xi_{T+k})\bigr]\\
 &\quad=\sum_{n\geq0}
      \mathbb E\bigl[\mathbf1_{A\cap\{T=n\}}
                    g(\xi_{n+1},\ldots,\xi_{n+k})\bigr]\\
 &\quad=\sum_{n\geq0}\mathbb P(A\cap\{T=n\})
                   \mathbb Eg(\xi_1,\ldots,\xi_k)
 =\mathbb P(A)\mathbb Eg(\xi_1,\ldots,\xi_k).
 \end{split}
 \]
The interchange is justified by boundedness of \(g\) and the
disjoint stopping events; the omitted infinite-time event is null.
This proves the claim for every finite cylinder of future
increments. The monotone-class extension to their generated
sigma-algebra proves it for the entire future sequence.

In a nearest-neighbor negative-drift walk, each successively
lower integer level is hit almost surely, as proved below.
The next passage duration is a measurable function of the
future increment sequence: its event of exceeding \(m\) is
determined by the first \(m\) future increments.
Translation invariance and the proved independence make each
duration have the one-level law, independently of all preceding
durations. Induction establishes the asserted joint independence.
\end{proof}
\end{annexitem}

\begin{annexitem}{annex:p15.os-finite-exit}{thm:p15-os-finite-exit}{Independent favorable blocks}
\begin{proof}
Let \(w=b-a\). From every interior lattice point, a block of
\(w\) consecutive \(+1\) increments forces exit at or before
the end of that block. Disjoint blocks are independent and each
has probability \(p^w>0\).
If the walk survives \(m\) blocks, none of those blocks can have
consisted entirely of \(+1\) increments. Thus
\(\mathbb P_x(\tau_{a,b}>mw)\leq(1-p^w)^m\).
The right side tends to zero, proving almost-sure exit.
For \(c>0\) as in the statement, partition the positive times
into blocks to bound
\[
 \mathbb E_xe^{c\tau_{a,b}}
 \leq 1+\sum_{m\geq0}e^{c(m+1)w}
                       \mathbb P_x(\tau_{a,b}>mw)
 \leq1+e^{cw}\sum_{m\geq0}
                 \{e^{cw}(1-p^w)\}^m<\infty.
 \]
Every positive integer power is bounded by a constant times
this exponential, so all the stated time moments are finite.
The deterministic endpoint-increment cases and a boundary start
follow directly from the path definition.
\end{proof}
\end{annexitem}

\begin{annexitem}{annex:p15.os-fair-exit}{thm:p15-os-fair-exit}{The two bounded-path sampling arguments}
\begin{proof}
The preceding exit bound already proves integrability of
\(\tau=\tau_{a,b}\). The stopped positions belong to \([a,b]\),
so they are uniformly bounded and converge to \(S_\tau\).
Bounded optional sampling of the linear martingale, then
bounded convergence, gives
\[
 x=\mathbb E_xS_\tau
   =a\,\mathbb P_x(S_\tau=a)+b\,\mathbb P_x(S_\tau=b).
 \]
The two exit probabilities sum to one, yielding the formulas.

The quadratic process \(S_n^2-n\) is a martingale, since
\(\mathbb E(S_{n+1}^2\mid\mathcal F_n)=S_n^2+1\).
At the bounded time \(\tau\wedge n\),
\[
 x^2=\mathbb E_xS_{\tau\wedge n}^2-\mathbb E_x(\tau\wedge n).
 \]
The squares are uniformly bounded by
\(\max(a^2,b^2)\), and the time terms increase to \(\tau\).
Bounded and monotone convergence give
\[
 \mathbb E_x\tau
 =a^2\frac{b-x}{b-a}+b^2\frac{x-a}{b-a}-x^2
 =(b-x)(x-a).
 \]
At \(x=0\), the first two terms simplify to \(-ab\),
retaining the source's particular calculation.

For \(S_0=0\) and a fixed positive integer \(b\),
\(\mathbb P_0(T_b<T_{-M})=M/(M+b)\to1\).
Each such event implies \(T_b<\infty\), so the latter has
probability one. Reflection proves the same for negative levels,
and translation gives every starting integer.
In the source's special case,
\(\mathbb P_1(T_M<T_0)=1/M\) whereas
\(\mathbb P_1(T_0<T_M)=(M-1)/M\).
These are complementary events, not two probabilities of the
same event.
\end{proof}
\end{annexitem}

\begin{annexitem}{annex:p15.os-fair-hitting}{thm:p15-os-fair-hitting}{The exponential root and binomial inversion}
\begin{proof}
The previous result proves \(T_1<\infty\) almost surely.
For \(\theta>0\), let
\(\varphi(\theta)=(e^\theta+e^{-\theta})/2=\cosh\theta>1\).
The exponential martingale
\(M_n=e^{\theta S_n}/\varphi(\theta)^n\) starts at one.
Before \(T_1\), the position is at most zero, and at \(T_1\)
it is one. Therefore \(0\leq M_{T_1\wedge n}\leq e^\theta\).
Bounded stopping and bounded convergence give
\[
 1=\mathbb EM_{T_1}
   =e^\theta\mathbb E\varphi(\theta)^{-T_1}.
 \]
The function \(\cosh\theta\) increases continuously from one
to infinity for \(\theta\geq0\).
For \(0<s<1\), choose its unique positive solution of
\(\varphi(\theta)=1/s\).
Writing \(u=e^\theta>1\), the equation is
\(su^2-2u+s=0\), whose admissible root is
\[
 u=\frac{1+\sqrt{1-s^2}}s.
 \]
The generating function is its reciprocal, not this root:
\[
 \mathbb Es^{T_1}=u^{-1}
                  =\frac{1-\sqrt{1-s^2}}s.
 \]

The proved binomial-series formula at exponent \(1/2\) gives,
for \(|s|<1\),
\[
 \frac{1-\sqrt{1-s^2}}s
 =\sum_{n\geq1}(-1)^{n-1}\binom{1/2}{n}s^{2n-1}.
 \]
The coefficient is
\[
 (-1)^{n-1}\binom{1/2}{n}
 =\frac{1\cdot3\cdots(2n-3)}{2^n n!}
 =\frac1{2n-1}\binom{2n}{n}2^{-2n},
 \]
where the empty product at \(n=1\) is one.
The expectation has the absolutely convergent probability series
\(\sum_{k\geq1}\mathbb P(T_1=k)s^k\).
Equality on \(0<s<1\) forces equality of coefficients:
if two convergent power series first differed at degree \(j\),
division by \(s^j\) and \(s\downarrow0\) would contradict
their equality. This proves the inversion and parity assertion.
At zero the function tends to zero; at one monotone convergence
gives \(\mathbb P(T_1<\infty)=1\).
If \(\mathbb ET_1\) were finite, Wald's first identity for the
centered walk would give \(\mathbb ES_{T_1}=0\), contradicting
\(S_{T_1}=1\). Hence its mean is infinite.
\end{proof}
\end{annexitem}

\begin{annexitem}{annex:p15.os-biased-exit}{thm:p15-os-biased-exit}{The exponential harmonic function and interval means}
\begin{proof}
For \(r=q/p\), independence gives
\[
 \mathbb E(r^{S_{n+1}}\mid\mathcal F_n)
 =r^{S_n}(pr+q/r)=r^{S_n}(q+p)=r^{S_n}.
 \]
Every finite-time value is integrable, since each position has
finite support. The stopped values at \(\tau=T_a\wedge T_b\)
are bounded between the two finite positive endpoint powers.
The exit bound proves \(\tau<\infty\).
Optional sampling and bounded convergence therefore give
\[
 r^x=\mathbb E_xr^{S_\tau}
             =r^a\alpha+r^b(1-\alpha).
 \]
Since \(r\ne1\) and \(a<b\), the denominator \(r^b-r^a\)
is nonzero; solving gives the lower-exit probability and its
complement. Wald's first identity is applicable because the
exit time is integrable. Thus
\(x+(p-q)\mathbb E_x\tau=a\alpha+b\beta\), which gives the mean.

For roulette \(r=(20/38)/(18/38)=10/9\).
The upper-exit probability from \(20\) between \(0,40\) is
\(\{r^{20}-1\}/\{r^{40}-1\}\).
Factor \(r^{40}-1=(r^{20}-1)(r^{20}+1)\) to get the
displayed exact simplified expression and its rounded decimal.
\end{proof}
\end{annexitem}

\begin{annexitem}{annex:p15.os-drift-hitting}{thm:p15-os-drift-hitting}{One-sided limits and an integrable global minimum}
\begin{proof}
Suppose \(p>q\), so \(0<r<1\), and start at zero.
For fixed \(a<0\), the events \(\{T_a<T_b\}\) increase as
integer \(b\to\infty\), and their union is \(\{T_a<\infty\}\).
Indeed the finite path up to a finite \(T_a\) has a finite maximum.
The interval formula and \(r^b\to0\) therefore give
\(\mathbb P_0(T_a<\infty)=r^{-a}\).
For fixed \(b>0\), sending \(a\to-\infty\) in the upper-exit
formula gives \(\mathbb P_0(T_b<\infty)=1\), because
\(r^a\to\infty\).
Nearest-neighbor steps imply
\(\{I\geq k\}=\{T_{-k}<\infty\}\), so
\[
 \mathbb EI=\sum_{k\geq1}\mathbb P(I\geq k)
                =\sum_{k\geq1}r^k=\frac r{1-r}.
 \]
In particular the global minimum is finite almost surely.

At the bounded time \(T_b\wedge n\), the linear martingale gives
\[
 \mathbb E_0S_{T_b\wedge n}
                       =(p-q)\mathbb E_0(T_b\wedge n).
 \]
The position lies between \(-I\) and \(b\).
This integrable bound and \(T_b<\infty\) justify dominated
convergence on the left, giving limit \(b\).
Monotone convergence on the right gives
\(b=(p-q)\mathbb E_0T_b\), and proves that expectation finite.
No hitting-time moment was assumed before this calculation.
Reflect the increments and translate the starting point to obtain
the negative-drift assertion and its mean \(x/(q-p)\).
\end{proof}
\end{annexitem}

\begin{annexitem}{annex:p15.os-downward-variance}{thm:p15-os-downward-variance}{Wald's second identity supplies the missing second moment}
\begin{proof}
Let \(\rho=q-p>0\). The preceding hitting result already gives
\(\mathbb E_xV_0=x/\rho<\infty\).
The independent centered increments
\(\eta_j=\xi_j+\rho\) have variance
\(v=1-\rho^2=4pq\). Their sum at \(V_0\) is
\[
 R_{V_0}=\sum_{j=1}^{V_0}(\xi_j+\rho)
                  =-x+\rho V_0.
 \]
Wald's second identity applies because the stopping time is
integrable, giving
\(\mathbb E_x(-x+\rho V_0)^2=v\,x/\rho<\infty\).
The bound
\(\rho^2V_0^2\leq2(-x+\rho V_0)^2+2x^2\)
first proves \(\mathbb E_xV_0^2<\infty\).
Expansion is now justified:
\[
 \rho^2\mathbb E_xV_0^2-2\rho x(x/\rho)+x^2
                                      =v x/\rho.
 \]
Subtracting the square of its mean gives
\(\operatorname{Var}_x(V_0)=vx/\rho^3\).
The source quadratic martingale
\((S_n+\rho n)^2-vn\), with initial value \(x^2\),
has the same stopped identity; the \(L^2\) argument above
justifies its limit.

The passage from \(x\) to zero consists of the successive
passages from \(j\) to \(j-1\), \(j=x,\ldots,1\).
All these times are finite, and the restart lemma makes their
durations independent and identically distributed.
Thus \(V_0\) is a sum of \(x\) copies of the one-level time.
Once its second moment has been established, variance additivity
also explains why the answer must be a constant times \(x\).
The cases \(p=0\) and \(x=0\) have deterministic times and
the stated zero variances.
\end{proof}
\end{annexitem}

\begin{annexitem}{annex:p15.os-downward-pgf}{thm:p15-os-downward-pgf}{The downward root and Catalan coefficients}
\begin{proof}
For \(\theta<0\), \(\varphi(\theta)>1\).
Indeed \(\varphi'(0)=p-q<0\) and strict convexity makes
\(\varphi'\) no larger than this negative value on
\((-\infty,0]\); hence \(\varphi\) decreases to one as
\(\theta\uparrow0\). Before \(V_0\), \(S_n\geq1\);
at \(V_0\), \(S_{V_0}=0\).
Thus the stopped exponential martingale
\[
 M_{V_0\wedge n}
       =e^{\theta S_{V_0\wedge n}}/
                          \varphi(\theta)^{V_0\wedge n}
 \]
is bounded between zero and one.
Its expectation is \(e^{\theta x}\), and dominated convergence
using almost-sure finiteness gives
\(e^{\theta x}=\mathbb E_x\varphi(\theta)^{-V_0}\).
At \(\theta=0\) both sides are one.

For \(0<s<1\), the equation
\(\varphi(\theta)=1/s\) with \(\theta<0\) has a unique solution.
For \(u=e^\theta\in(0,1)\), it is
\(psu^2-u+qs=0\).
The admissible root is
\[
 u=\frac{1-\sqrt{1-4pqs^2}}{2ps},
 \]
as can be seen either from \(u\to0\) at \(s\downarrow0\),
or from the monotone correspondence just established.
This gives \(G_{V_0}(s)=u^x=f(s)^x\).

For \(|s|<1\), the binomial series is absolutely convergent,
since \(4pq<1\), and its coefficients give
\[
 f(s)=\sum_{k\geq0} C_kp^kq^{k+1}s^{2k+1}.
 \]
To check the coefficient, use
\[
 \frac{(-1)^k}{2p}\binom{1/2}{k+1}(4pq)^{k+1}
 =\frac{(2k)!}{k!(k+1)!}p^kq^{k+1}.
 \]
Power-series uniqueness gives the one-level probabilities.
The restart lemma writes the time from \(x\) as a sum of \(x\)
independent one-level times; multiplying their absolutely
convergent series gives the displayed finite convolution
coefficient at \(x+2k\), and zero at other times.
It also explains structurally the power \(f(s)^x\).
The endpoint values follow from \(V_0\geq1\) at \(x\geq1\)
and its almost-sure finiteness. Initial position zero and the
deterministic case \(p=0\) follow from the path itself.
\end{proof}
\end{annexitem}

\begin{annexitem}{annex:p15.os-quartic-time}{thm:p15-os-quartic-time}{The coefficient calculation and justified unbounded limit}
\begin{proof}
For a fair sign independent of the past,
\[
 \mathbb E(S_{n+1}^2\mid\mathcal F_n)=S_n^2+1,\qquad
 \mathbb E(S_{n+1}^4\mid\mathcal F_n)=S_n^4+6S_n^2+1.
 \]
All finite-time moments exist. For
\(Y_n=S_n^4-6nS_n^2+bn^2+cn\), subtraction of \(Y_n\)
from its next conditional expectation gives
\((2b-6)n+(b+c-5)\).
It vanishes for every \(n\) exactly when \(b=3,c=2\).

The finite-interval exit bound proves every moment of \(T\)
finite before any unbounded optional identity is used.
The fair-exit formula gives \(\mathbb ET=a^2\), and
\(|S_{T\wedge n}|\leq a\).
At the bounded time \(T\wedge n\), \(\mathbb EY_{T\wedge n}=0\).
Its terms are dominated respectively by
\(a^4,6a^2T,3T^2,2T\), all integrable.
Dominated convergence therefore gives
\[
 0=a^4-6a^2\mathbb ET+3\mathbb ET^2+2\mathbb ET
   =-5a^4+3\mathbb ET^2+2a^2.
 \]
Solving gives the second moment; subtracting \(a^4\)
gives the variance. At \(a=1\), \(T=1\) deterministically,
in agreement with both formulas.
\end{proof}
\end{annexitem}

\begin{annexitem}{annex:p15.os-exponential-exit}{thm:p15-os-exponential-exit}{Integrable exponential overshoots and infinite-time ruin}
\begin{proof}
Write \(\theta=\theta_0<0\).
The equation \(\mathbb Ee^{\theta\xi_1}=1\), with a
nonconstant increment, implies positive probability of both
strictly positive and strictly negative increments.
For if \(\xi_1\geq0\) almost surely, the exponential is at
most one and equality of its mean to one would force
\(\xi_1=0\) almost surely; the nonpositive case is analogous.
Choose \(\delta>0\) with
\(p_\delta=\mathbb P(\xi_1\geq\delta)>0\).
A fixed number \(k\) of such steps, with \(k\delta>b-a\),
forces exit from the interval. Independent disjoint blocks give
\[
 \mathbb P_x(\tau>mk)\leq(1-p_\delta^k)^m,
 \]
so \(\mathbb E_x\tau<\infty\).

The exponential process \(X_n=e^{\theta S_n}\) is integrable,
since independence and the nonnegative product integral give
\(\mathbb E_xX_n=e^{\theta x}\varphi(\theta)^n=e^{\theta x}\).
Independence of the next increment gives
\(\mathbb E(X_{n+1}\mid\mathcal F_n)=X_n\varphi(\theta)=X_n\).
Thus it is a nonnegative martingale without requiring
\(\mathbb E|\xi_1|<\infty\).
On \(\{\tau>n\}\), \(S_n>a\) and
\(X_n\leq e^{\theta a}\); thus the residual family
\(X_n\mathbf1_{\{\tau>n\}}\) is bounded.
The terminal overshoot need not be bounded, so we prove
its integrability separately. For \(n\geq1\),
\[
 \begin{split}
 \mathbb E_x(e^{\theta S_n};\tau=n)
 &\leq\mathbb E_x(e^{\theta S_n};\tau>n-1)\\
 &=\mathbb E_x(e^{\theta S_{n-1}};\tau>n-1)
 \leq e^{\theta a}\mathbb P_x(\tau>n-1).
 \end{split}
 \]
The equality uses independence of the next increment and
\(\varphi(\theta)=1\).
Summing nonnegative terms gives
\(\mathbb E_xe^{\theta S_\tau}\leq e^{\theta a}\mathbb E_x\tau<\infty\).
The residual-tail stopping theorem now yields the exact identity
\(\mathbb E_xe^{\theta S_\tau}=e^{\theta x}\).
On the lower-exit event,
\(e^{\theta S_\tau}\geq e^{\theta a}\), so its probability
is bounded as stated.

For the possibly infinite lower-crossing time \(T\), the
nonnegative martingale has a finite almost-sure terminal limit
by the nonnegative supermartingale theorem.
Its sampling inequality gives
\[
 e^{\theta a}\mathbb P_x(T<\infty)
 \leq\mathbb E_x X_T\leq e^{\theta x}.
 \]
The left side counts only finite crossing outcomes; the terminal
limit on \(\{T=\infty\}\) is nonnegative and need not be zero.
Thus the argument does not require finite mean time or bounded
overshoots for this one-sided bound.
\end{proof}
\end{annexitem}

\begin{annexitem}{annex:p15.os-skip-free-ruin}{thm:p15-os-skip-free-ruin}{Existence of the root and the surviving-path limit}
\begin{proof}
Because \(\xi_1\geq-1\), \(\varphi(\theta)\) is finite
for every \(\theta\leq0\). It is continuous there, by
dominated convergence on each compact parameter interval.
At negative arguments it is twice differentiable: the
positive-tail derivatives are dominated on compact subintervals
of \((-\infty,0)\) by constants times \(x^2e^{-cx}\),
bounded for \(x\geq0\); the negative part has only the value
\(-1\). Hence
\(\varphi''(\theta)=\mathbb E(\xi_1^2e^{\theta\xi_1})>0\),
so it is strictly convex there.
Also \(\varphi(\theta)\to\infty\) as \(\theta\to-\infty\)
because \(\mathbb P(\xi_1=-1)>0\).

For \(s>0\),
\[
 \frac{1-\varphi(-s)}s
 =\mathbb E\frac{1-e^{-s\xi_1}}s.
 \]
On nonnegative increments the integrand increases to \(\xi_1\)
as \(s\downarrow0\), by
\((1-e^{-sx})/s=\int_0^x e^{-su}\,du\).
On the only negative value, it tends to \(-1\).
Monotone convergence on the positive part gives limit
\(\mathbb E\xi_1\in(0,\infty]\).
Therefore \(\varphi(-s)<1\) for sufficiently small \(s>0\).
Continuity supplies a negative root of \(\varphi(\theta)=1\);
strict convexity, together with \(\varphi(0)=1\), makes it
unique among negative arguments.

The strong law, with its extended-mean version when necessary,
gives \(S_n\to+\infty\) almost surely: the negative part of
\(\xi_1\) is bounded and its mean is strictly positive.
For integers \(a<x\), a downward step of at most one means
that the first crossing of \(a\) is exactly a hit of \(a\).
Before \(T_a\) the walk is above \(a\); thus
\[
 0\leq e^{\theta_0S_{T_a\wedge n}}\leq e^{\theta_0a}.
 \]
On \(\{T_a<\infty\}\) this tends to \(e^{\theta_0a}\).
On \(\{T_a=\infty\}\), the positive-drift limit makes it tend
to zero. Bounded stopping and dominated convergence give
\[
 e^{\theta_0x}
     =e^{\theta_0a}\mathbb P_x(T_a<\infty).
 \]

For the three-payoff game, put \(r=e^{\theta_0}\in(0,1)\).
The root equation is
\((r^{-1}+r+r^2)/3=1\), or
\[
 r^3+r^2-3r+1=(r-1)(r^2+2r-1)=0.
 \]
The unique root in \((0,1)\) is \(\sqrt2-1\).
The increment mean is \(2/3>0\), so all hypotheses hold
and the formula from capital \(i\) follows.
\end{proof}
\end{annexitem}

\begin{annexitem}{annex:p15.os-insurance}{thm:p15-os-insurance}{The normal exponential moment and the ruin event}
\begin{proof}
The increments \(\xi_j=c-\zeta_j\) are i.i.d. normal with
mean \(d=c-\mu>0\) and variance \(\sigma^2>0\).
Completing the square in their density gives, for every real
\(\theta\),
\[
 \begin{split}
 \mathbb Ee^{\theta\xi_1}
 &=\frac1{\sqrt{2\pi}\sigma}
   \int_{\mathbb R}
     \exp\left(\theta u-\frac{(u-d)^2}{2\sigma^2}\right)\,du\\
 &=e^{d\theta+\sigma^2\theta^2/2}
   \frac1{\sqrt{2\pi}\sigma}
   \int_{\mathbb R}
     \exp\left(-\frac{(u-d-\sigma^2\theta)^2}{2\sigma^2}\right)\,du\\
 &=e^{d\theta+\sigma^2\theta^2/2}.
 \end{split}
 \]
The last integral is the proved normal normalization after an
affine one-variable substitution.
The negative root making this moment one is
\(\theta_0=-2d/\sigma^2\).
The exponential lower-crossing bound with level \(a=0\)
and initial assets \(x>0\) therefore gives
\(\mathbb P_x(\text{ruin})\leq e^{\theta_0x}\).
This treats ruin as reaching zero or any negative value,
not merely hitting zero exactly.
If the variance is zero, every increment equals \(d>0\),
so the deterministic assets cannot be ruined.
\end{proof}
\end{annexitem}

\begin{annexitem}{annex:p15.os-boundaries}{thm:p15-os-boundaries}{Four explicit failures of omitted hypotheses}
\begin{proof}
The first example follows from the proved fair hitting law:
\(T_1<\infty\), \(S_{T_1}=1\), and \(\mathbb ET_1=\infty\).
If its stopped walk were uniformly integrable, the sampling
theorem would give terminal mean zero, a contradiction.
Thus its residual family cannot meet the residual-tail criterion.

For the walk from one stopped at zero, bounded stopping gives
\(\mathbb ES_{T_0\wedge n}=1\).
It remains nonnegative and tends almost surely to zero, since
the fair hitting theorem gives \(T_0<\infty\).
Its terminal mean is strictly smaller than its initial mean,
as permitted by nonnegative supermartingale sampling.

For doubling, let the independent fair signs be successes and
failures and put \(T=\inf\{j\geq1:\xi_j=1\}\).
At step \(j\), wager the predictable amount
\(2^{j-1}\mathbf1_{\{T\geq j\}}\).
The cumulative gain at horizon \(n\) is
\[
 G_n=
 \begin{cases}
 1,&T\leq n,\\
 -(2^n-1),&T>n.
 \end{cases}
 \]
It is a martingale, since each finite predictable wager
multiplies a centered independent sign.
Also \(\mathbb P(T>n)=2^{-n}\), so
\(\mathbb ET=\sum_{n\geq0}2^{-n}=2\),
\(\mathbb EG_n=0\), and \(G_T=1\).
On the event of survival to step \(j\), the conditional mean
absolute increment is \(2^{j-1}\), so no deterministic finite
bound uniform in \(j\) exists.
Indeed \(\mathbb E|G_n-1|=2^n\mathbb P(T>n)=1\),
which also explicitly excludes \(L^1\) convergence and uniform
integrability of these stopped gains.

Finally \(\{T=0\}=\{\xi=-1\}\) in the future-selected example
has probability \(1/2\) and is not in the trivial
\(\mathcal F_0\). Thus that bounded random time is not a stopping
time. Its sampled value is one on \(\{\xi=1\}\) and zero
otherwise, giving mean \(1/2\).
\end{proof}
\end{annexitem}
